% Define a macro to combine section number and title with background color
\newcommand{\subsectionformat}[1]{%
  \colorbox{LightBlue!40}{\parbox{\dimexpr\textwidth-2\fboxsep}{#1}}%
}

% Apply the macro in \titleformat without using # directly
\titleformat{\subsection}[block]
  {\bfseries} % Font style: bold italic
  {} % Leave label empty in \titleformat
  {-1pt}
  {\subsectionformat} % Call the combined format with background
\titlespacing*{\subsection}{0pt}{-1pt}{-1pt}

% Define a macro to combine section number and title with background color
% \newcommand{\sectionformat}[1]{%
%   \colorbox{LightBlue!90}{\parbox{\dimexpr\textwidth-2\fboxsep}{\thesection\ #1}}%
% }

\newif\ifsectionnumbered
\sectionnumberedtrue % default: show number

\newcommand{\sectionformat}[1]{%
  \colorbox{LightBlue!90}{%
    \parbox{\dimexpr\textwidth-2\fboxsep}{%
      \ifsectionnumbered
        \thesection\ #1%
      \else
        #1%
      \fi
    }%
  }%
}

% Then define a clean unnumbered command:
\newcommand{\unnumberedsection}[1]{%
  \sectionnumberedfalse
  \section*{#1}%
  \sectionnumberedtrue
}

% Apply the macro in \titleformat without using # directly
\titleformat{\section}[block]
  {\bfseries} % Font style: bold italic
  {} % Leave label empty in \titleformat
  {0pt}
  {\sectionformat} % Call the combined format with background
\titlespacing*{\section}{0pt}{0pt}{0pt}

\newcommand{\unnumberedsection}[1]{%
  \vspace{...}% match your spacing
  {\bfseries\sectionformat{#1}}%
  \addcontentsline{toc}{section}{#1}% optional: add to TOC
}
                  % Title format
